\documentclass[10pt,a4paper]{amsart}
\usepackage[margin=19mm]{geometry}
\usepackage{amsmath,amssymb,mathtools}
\usepackage[scr=rsfs,cal=euler]{mathalfa}
\usepackage{xfrac}
\usepackage[unicode,hidelinks,bookmarks=false]{hyperref}
\newcommand{\ie}{i.e.\ }
\newcommand{\cf}{cf.\ }
\newcommand{\resp}{respectively\ }
\newcommand{\Subsection}{\subsection}
\providecommand{\nobreakdash}{\nobreak\hbox{-}}
\setlength{\emergencystretch}{2em}
\multlinegap0pt
\usepackage{amssymb}
\usepackage{mathtools}
\usepackage{xcolor}
\usepackage{cancel}
\usepackage{float}
\usepackage[shortlabels]{enumitem}
\newtheorem{thm}{Theorem}
\newtheorem{lemma}{Lemma}
\newcommand {\Div}  { {\rm div} }
\renewcommand{\color}[2][]{}
\newcommand{\eps}{\varepsilon}
\title{Sobolev estimates for the Keller-Segel system and applications to the JKO scheme}
\author{Charles Elbar}
\date{Source: arXiv:2410.15095v3 (CC BY 4.0)}
\begin{document}
\maketitle

\noindent\emph{Source and licence.} Sobolev estimates for the Keller-Segel system and applications to the JKO scheme. Version arXiv:2410.15095v3 (CC BY 4.0), distributed by arXiv under the Creative Commons Attribution 4.0 International licence, http://creativecommons.org/licenses/by/4.0/, as recorded in the arXiv licence metadata for this exact version and shown on the abstract page https://arxiv.org/abs/2410.15095v3. Full licence notice: \texttt{licenses/arxiv-2410.15095-CC-BY-4.0.txt}.\par\smallskip

\noindent\textit{Editorial scope.} Counted unit: the article's lemma lem:strict\_positivy (source0.tex lines 1039--1043). The article's complete original proof of it is delivered with it (source0.tex lines 1045--1203, one \texttt{proof} environment of 159 lines). No mathematics is written, repaired or re-derived: the statement and the proof are the article's own lines, in the article's own notation, and no retained line is substituted or re-flowed.  Retained uncounted as the source-local context the proof needs: lines 213--215; lines 443--445; lines 448--455; lines 963--965; lines 984--986.  Nothing was omitted from any retained span, so the package carries no omission record.  No retained line contains a citation, so the wrapper carries no bibliography and no bibliography entry.  No retained line contains a figure environment, an image-inclusion command or drawing code, so no image is delivered and the package declares no asset dependency.  Editorial formatting only, in addition: an AMS \texttt{amsart} preamble that restates the article's own declarations and macros used by the retained lines, namely the environments \texttt{thm}, \texttt{lemma} on separate counters, together with the standard packages those lines need.  The retained environments print in this document as Theorem 1, Lemma 1 in order of appearance, whereas the article prints the same items with the numbering of its own full document.  The complete native source of the article as distributed in the arXiv e-print for this exact version is bundled unchanged as \texttt{sources/167101/source0.tex}.
\begin{equation}\label{eq:KS}
\partial_t \rho - \Delta \rho + \Div(\rho\nabla u[\rho]) = 0, \quad -\Delta u[\rho] =\rho.
\end{equation}

\begin{equation}\label{eq:def_JKO2}
\rho_\tau(t) = \rho_\tau^{n+1} \text{ for } t \in (n\tau, (n+1)\tau],
\end{equation}

\begin{thm}[JKO scheme]\label{thm:JKO}
Let $a\le \rho_{0}\le b$, where $a,b>0$ be an initial condition in $L^1(\Omega)$ and such that $\mathcal{F}_{p}[\rho(0)]<+\infty$ for all $1\le p<+\infty$. {\color{blue} Let $\rho_{\tau}$ the constructed curve~\eqref{eq:def_JKO2} of the JKO scheme.   Then there exists  $M_d,\nu >0$ in the definition of $\widetilde{\mathcal{J}}$  and $\eps_{d}$ depending on the dimension and the final time of existence $T$ such that if $\|\rho_0\|_{L^{\frac{d}{2}}(\Omega)}\le \eps_d$:
\begin{itemize}
    \item $\|\rho_{\tau}^{n}\|_{L^{\frac d2}}^2\le \frac{M_d}{C_0}$ for some $C$ independent of $n$ and $\widetilde{\mathcal J}[\rho_n]=\mathcal J[\rho_n]$ so the artificial term disappears
    \item $\rho_{\tau}$ is bounded uniformly in $\tau$ in $L^{\infty}((0,T)\times\Omega)\cap L^{\infty}(0,T; W^{1,p}(\Omega))\cap L^{2}(0,T; H^{2}(\Omega))$ for all $1\le p<+\infty$
    \item $\rho_{\tau}$ converges to $\rho$ as $\tau\to 0$ strongly in $L^{2}(0,T; H^{2}(\Omega))$, where $\rho$ is a weak solution of the KS system~\eqref{eq:KS}.
\end{itemize}}
\end{thm}

\begin{equation}\label{eq:JKO_Ld2_1}
 \|\rho_k\|_{L^{\frac d2}}^2\le M_d/C_0   .
\end{equation}

\begin{equation}\label{eq:JKO_Ld2_2}
\|\rho_k\|_{L^{\frac d2}}^2\le C\nu+ C\nu \widetilde{\mathcal{J}}[\rho_\tau^n]+ M_d.
\end{equation}

\begin{lemma}[Strict positivity of the JKO scheme]\label{lem:strict_positivy}
Let $\rho^{n+1}_{\tau}\in Prox_{\widetilde{\mathcal{J}}}^{\tau}(\rho^{n}_{\tau})$ be the penalized JKO sequence
associated to the Keller-Segel system, with initial density $\rho_0$ as in Theorem~\ref{thm:JKO}.
Then $\rho^n_\tau>0$ almost everywhere in $\Omega$ for every $n\ge 0$.
\end{lemma}

\begin{proof}
Fix $n\ge 0$ and denote
\[
\mu:=\rho_\tau^n,
\qquad
\rho:=\rho_\tau^{n+1}\in Prox_{\widetilde{\mathcal J}}^\tau(\mu).
\]
That is, $\rho$ minimizes the penalized JKO scheme
\begin{equation}\label{eq:def_JKO_penalized}
\widetilde{\mathcal J}_\tau(\sigma\,|\,\mu)
:=\frac{1}{2\tau}W_2^2(\sigma,\mu)
+\mathcal J[\sigma]
+\frac1\nu\bigl(C_0\|\sigma\|_{L^{d/2}(\Omega)}-M_d\bigr)_+,
\end{equation}
among all densities $\sigma\ge 0$ with $\int_\Omega \sigma\,dx=M$, where $M$ is the mass of $\rho_0$ (and $\mu$).

We argue by contradiction. Assume that the zero set
\[
A:=\{x\in\Omega:\rho(x)=0\}
\]
has positive measure, $|A|>0$.

Let $\eta\equiv M/|\Omega|$ be the constant density with mass $M$.
For $\varepsilon\in(0,1)$ define the convex perturbation
\[
\rho_\varepsilon:=(1-\varepsilon)\rho+\varepsilon\eta.
\]
Then $\rho_\varepsilon\ge \varepsilon\eta>0$ a.e., $\rho_\varepsilon$ has the same mass $M$, and therefore it is an admissible
competitor in~\eqref{eq:def_JKO_penalized}. By minimality of $\rho$,
\begin{equation}\label{eq:minimality_start}
0\le \widetilde{\mathcal J}_\tau(\rho_\varepsilon\,|\,\mu)-\widetilde{\mathcal J}_\tau(\rho\,|\,\mu).
\end{equation}
We show that the right-hand side is strictly negative for $\varepsilon>0$ small enough, yielding a contradiction.

\medskip
\underline{Step 1: Wasserstein term.}
We use convexity of $W_2^2(\,\cdot\,,\mu)$ with respect to the first argument.
Let $\gamma_\rho$ be an optimal transport plan between $\rho$ and $\mu$, and $\gamma_\eta$ an optimal plan between $\eta$ and $\mu$.
Then
\[
\gamma_\varepsilon:=(1-\varepsilon)\gamma_\rho+\varepsilon\gamma_\eta
\]
is a transport plan between $\rho_\varepsilon$ and $\mu$, hence
\begin{align*}
W_2^2(\rho_\varepsilon,\mu)
&\le \int_{\Omega\times\Omega}|x-y|^2\,d\gamma_\varepsilon(x,y)\\
&=(1-\varepsilon)\int |x-y|^2\,d\gamma_\rho+\varepsilon\int |x-y|^2\,d\gamma_\eta\\
&=(1-\varepsilon)W_2^2(\rho,\mu)+\varepsilon W_2^2(\eta,\mu).
\end{align*}
Therefore
\begin{equation}\label{eq:W2_Oeps}
\frac{1}{2\tau}\Bigl(W_2^2(\rho_\varepsilon,\mu)-W_2^2(\rho,\mu)\Bigr)
\le \frac{\varepsilon}{2\tau}\,W_2^2(\eta,\mu)
=C\,\varepsilon,
\end{equation}


\medskip
\underline{Step 2: Aggregation term.}
Denote the Keller--Segel interaction energy by
\[
I(\sigma):=-\frac12\int_\Omega \sigma\,u[\sigma]\,dx,
\]
Expanding,
\begin{align*}
I(\rho_\varepsilon)=(1-\varepsilon)^2 I(\rho)
-\frac{\varepsilon(1-\varepsilon)}2\int_\Omega \rho\,u[\eta]\,dx
-\frac{\varepsilon(1-\varepsilon)}2\int_\Omega \eta\,u[\rho]\,dx
+\varepsilon^2 I(\eta).
\end{align*}
Define the cross term
\[
B(\rho,\eta):=-\frac12\int_\Omega \rho\,u[\eta]\,dx=-\frac12\int_\Omega \eta\,u[\rho]\,dx.
\]
Then we obtain

\begin{align}
I(\rho_\varepsilon)-I(\rho)=2\varepsilon\bigl(B(\rho,\eta)-I(\rho)\bigr)+O(\varepsilon^2).\label{eq:I_diff_Oeps}
\end{align}
Note that~\eqref{eq:JKO_Ld2_1}-\eqref{eq:JKO_Ld2_2} yield $L^{d/2}(\Omega)$ bounds on $\rho$, and $\eta$ is bounded in $L^{d/2}(\Omega)$ by definition. Thus
\[
|I(\rho)|\le C\|\rho\|_{L^{d/2}}^2,\qquad |B(\rho,\eta)|\le C\|\rho\|_{L^{d/2}}\|\eta\|_{L^{d/2}},
\]
so the right-hand side of~\eqref{eq:I_diff_Oeps} is bounded by $C\varepsilon$ for a constant $C$ independent of $\varepsilon$:
\begin{equation}\label{eq:I_Oeps}
I(\rho_\varepsilon)-I(\rho)\le C\varepsilon.
\end{equation}

\medskip
\underline{Step 3: Penalization term.}
Define the penalty functional
\[
P(\sigma):=\frac{1}{\nu}\bigl(C_0\|\sigma\|_{L^{d/2}(\Omega)}-M_d\bigr)_+.
\]
Then
\begin{align}
P(\rho_\varepsilon)-P(\rho)
&\le C\|\rho_\varepsilon-\rho\|_{L^{d/2}}
= C\|\varepsilon(\eta-\rho)\|_{L^{d/2}}
= C\,\varepsilon.
\label{eq:penalty_Oeps}
\end{align}
\medskip
\underline{Step 4: Entropy.}
Let $f(s)=s\log s$ with the convention $f(0)=0$.
We estimate $\int_\Omega (f(\rho_\varepsilon)-f(\rho))$ by splitting $\Omega=A\cup(\Omega\setminus A)$.

\smallskip
On $A$:
Since $\rho=0$ on $A$, we have $\rho_\varepsilon=\varepsilon\eta$ on $A$, hence
\begin{align}
\int_A f(\rho_\varepsilon)\,dx
&=\int_A \varepsilon\eta\log(\varepsilon\eta)\,dx
=\varepsilon\log\varepsilon\int_A \eta\,dx+\varepsilon\int_A \eta\log\eta\,dx.
\label{eq:entropy_on_A}
\end{align}

On $\Omega\setminus A$:
By convexity of $f$ on $[0,\infty)$,
\[
f\bigl((1-\varepsilon)\rho+\varepsilon\eta\bigr)\le (1-\varepsilon)f(\rho)+\varepsilon f(\eta),
\]
hence
\[
f(\rho_\varepsilon)-f(\rho)\le \varepsilon\bigl(f(\eta)-f(\rho)\bigr).
\]
Integrating over $\Omega\setminus A$ gives
\begin{equation}\label{eq:entropy_off_A}
\int_{\Omega\setminus A}\bigl(f(\rho_\varepsilon)-f(\rho)\bigr)\,dx
\le \varepsilon\int_{\Omega\setminus A}\bigl(f(\eta)-f(\rho)\bigr)\,dx
\le C\,\varepsilon,
\end{equation}
the last estimate following from~\eqref{eq:JKO_Ld2_1}-\eqref{eq:JKO_Ld2_2}.

Combining~\eqref{eq:entropy_on_A} and~\eqref{eq:entropy_off_A}, we obtain
\begin{equation}\label{eq:entropy_total}
\int_\Omega \rho_\varepsilon\log\rho_\varepsilon-\rho\log\rho\,dx
\le
\varepsilon\log\varepsilon\int_A \eta\,dx + C\,\varepsilon,
\end{equation}
for some constant $C$ independent of $\varepsilon$.

\medskip
\underline{Step 5: Contradiction.}
Putting everything together, we obtain
\[
0\le
\widetilde{\mathcal J}_\tau(\rho_\varepsilon\,|\,\mu)-\widetilde{\mathcal J}_\tau(\rho\,|\,\mu)
\le
\varepsilon\log\varepsilon\int_A \eta\,dx + C\varepsilon,
\]
for some constant $C$ independent of $\varepsilon$. Since $\eta\equiv M/|\Omega|$, we have
\[
\int_A \eta\,dx=\frac{M}{|\Omega|}|A|>0.
\]
Hence, there exists $\varepsilon_0\in(0,1)$ such that the right-hand side is strictly negative for all
$\varepsilon\in(0,\varepsilon_0)$, contradicting~\eqref{eq:minimality_start}.
Therefore $|A|=0$, i.e.\ $\rho=\rho_\tau^{n+1}>0$ almost everywhere.
\end{proof}

\end{document}
